Continuous games are multiplayer games in which strategy sets are compact and utility functions are continuous. Many application domains have a~continuum of actions, such as the allocation of time, resources, location in space \cite{Kamra18}, or parameters of classifiers \cite{Loiseau19}. This also includes several games modeling cybersecurity scenarios \cite{Niu2021,Loiseau19,roussillon2021scalable}.

Continuous games have equilibria in mixed strategies by Glicksberg's generalization of Nash's theorem \cite{Glicksberg52}, but these equilibria are typically very hard to characterize and compute. We point out the main difficulties in the development of algorithms and numerical methods for continuous games:
\begin{itemize}
  \item The equilibrium can be any tuple of mixed strategies with infinite supports~\cite{Karlin59} or almost any tuple of~finitely-supported mixed strategies \cite{Rehbeck18}.
  \item Bounds on the size of supports of equilibrium strategies are known only for particular classes of continuous games \cite{stein.ozdaglar.ea:separable}.
  \item Some important games have only mixed equilibria; for example, certain variants of Colonel Blotto games \cite{Golman2009}.
  \item Finding a mixed strategy equilibrium involves locating its support, which lies inside a continuum of points.
\end{itemize}

To the best of our knowledge, algorithms and numerical methods exist only for very special classes of continuous games. In particular, two-player zero-sum polynomial games can be solved using sum-of-squares optimization based on a sequence of semidefinite relaxations of the original problem \cite{parrilo:polynomial,LarakiLasserre12}. \textcite{BasarOlsder99} have written a detailed analysis of equilibria for some families of games with a particular shape of utility functions (e.g., games of timing, bell-shaped kernels). Fictitious play, one of the principal learning methods for finite games, was recently extended to continuous action spaces and applied to Blotto games~\cite{Ganzfried21}. The dynamics of fictitious play have been analyzed only under further restrictive assumptions in the continuous setting \cite{Hofbauer06}. No-regret learning, studied by \textcite{Mertikopoulos2019}, can be applied to finding pure equilibria or to mixed strategy learning in finite games. Convergence guarantees for algorithms in the distributed environment solving convex-concave games and some generalizations thereof have been developed by \textcite{mertikopoulos.lecouat.ea:optimistic}. In a similar setting, \textcite{chasnov.ratliff.ea:convergence} provides convergence guarantees to a neighborhood of a stable Nash equilibrium for gradient-based learning algorithms.
